%% Problem formulation, the three dimensions of an observation process, and the re-observation protocol.
%% Included by sigconf.tex. Do not add a preamble here.

\section{Observation-Process Shift}
\label{sec:problem}

\subsection{Problem formulation}
Let $X=\{x(t)\}_{t\in\mathcal{T}}$, $x(t)\in\mathbb{R}^{N}$, be an underlying multivariate temporal process over
a window $\mathcal{T}$. The learner never observes $X$. It receives a finite set of timestamped, partially
observed measurements
\begin{equation}
\Oobs=\{(t_k,\,x(t_k),\,m_k)\}_{k=1}^{K},\qquad m_k\in\{0,1\}^{N},
\end{equation}
where $m_{k,i}=1$ iff variable $i$ was measured at time $t_k$. The mechanism that selects $\{(t_k,m_k)\}$ is the
\emph{observation process} $\Sobs$; we write $\Oobs=\Sobs(X)$ and model it as a random variable
$\Sobs\sim P(\Sobs)$. A model $f_\theta$ maps $\Oobs$ to a task output $\hat Y$, either a forecast of $x$ over a
future window or a label.

Training and test data are drawn from joint distributions over $(X,\Sobs)$, and we separate two ways they can
differ:
\begin{equation}
P_{\mathrm{train}}(X)\ \neq\ P_{\mathrm{test}}(X)
\qquad\textit{(process shift)},
\end{equation}
\begin{equation}
P_{\mathrm{train}}(\Sobs)\ \neq\ P_{\mathrm{test}}(\Sobs)
\qquad\textit{(observation-process shift)}.
\end{equation}
This paper concerns the second, with $P(X)$ comparable. A protocol that draws training and test data from one
observation process cannot see it.

\subsection{Three dimensions of an observation process}
\label{sec:family}
An observation process is not a scalar. Because $\Sobs$ acts on a mask indexed by both time and variable, it can
differ from another process in three elementary ways: in \emph{how much} it records, in \emph{when} it records,
and in \emph{which variables} it records. We organise both the process family and the evaluation protocol around
those three dimensions, so that a measured degradation can be attributed to one of them.

We require the processes to be \emph{realisable}: each must correspond to a way the same trajectory could have
been recorded, so that the target stays well defined. We therefore admit only processes that act on the mask by
\emph{removal},
\begin{equation}
\Sobs:(t,m)\ \mapsto\ (t,\,m\odot\kappa),\qquad \kappa\in\{0,1\}^{K\times N},
\end{equation}
with values at removed entries discarded rather than imputed. Every $\Sobs(X)$ is then a sub-sample of the
original record and no artificial value is introduced, which is what permits one held-out set to be re-observed
many times with its supervision untouched.

Nine primitives instantiate the three dimensions:
\begin{itemize}
\item \textbf{density} --- how much survives. \textsc{RandomDropout}$(p)$ removes each observation independently
with probability $p$; \textsc{DensityControlledSampling}$(\rho)$ retains a fixed fraction $\rho$ of every
variable's observations, fixing density independently of the original pattern.
\item \textbf{timing} --- when observations arrive. \textsc{TemporalStride}$(k)$ keeps every $k$-th observation
of each variable, lowering the sampling \emph{rate} while keeping it regular; \textsc{BurstMissingness}$(p)$
removes contiguous windows covering a fraction $p$ of the observed span; \textsc{LongGapMissingness}$(p)$ is the
single-window case; \textsc{LocalTemporalMask}$(p)$ removes one contiguous run per variable, matching random
dropout in density but clustering the loss within each channel.
\item \textbf{variable availability} --- which channels exist at all. \textsc{VariableDropout}$(p)$ removes whole
variables independently; \textsc{VariableGroupDropout}$(p)$ removes contiguous blocks of variables together,
modelling a failed monitor or an unavailable panel.
\end{itemize}
Processes compose sequentially ($\Sobs_2\circ\Sobs_1$, written \texttt{a+b}) and mix
($\Sobs\sim\mathrm{Unif}\{\Sobs_1,\dots\}$, written \texttt{a|b}), so an observation regime is specified by a
string rather than by code, and a process can move two dimensions at once.

\paragraph{The dimensions are not separable at the level of a single variable.}
It is tempting to summarise what a process did to variable $i$ by how many of its observations survived, and to
expect the error on variable $i$ to follow. It does not, because $f_\theta$ predicts each variable from the whole
observation. Two processes can leave one channel exactly the same observations and still produce very different
error on it, if they differ in what they removed from the \emph{other} channels
(Figure~\ref{fig:views}). A per-variable missingness rate is therefore not a sufficient statistic for the damage,
which is the concrete sense in which the observation process has to be treated as a structured object.

\subsection{Controlled re-observation and validity}
\label{sec:reobs}
Observation-process shift is measured by re-observation. For an evaluation process $\Sobs$ we replace the
\emph{history} of each held-out input by $\Sobs$ applied to it, and leave the prediction times, forecast targets
and labels exactly as they were. The comparison therefore isolates the observation process: the underlying
trajectory and the supervision are held fixed by construction, and a single test set yields as many shifted
conditions as there are processes. The pseudo-random draw is seeded independently of the run, so every model is
scored on byte-identical re-observed inputs and comparisons are paired at the level of individual test inputs.
The \emph{clean} condition is the dataset's own observation process, $\Sobs=\mathrm{id}$.

Because two of our mechanisms select observation processes rather than merely sampling them, we also need to rule
out degenerate choices. We impose validity constraints $\mathcal{C}$ on relative density, variable coverage,
maximum gap and observation count,
\begin{equation}
\mathcal{C}:\quad \rho_{\min}\le\frac{\rho(\Sobs(X))}{\rho(X)}\le\rho_{\max},\quad
\Delta t_{\max}(\Sobs(X))\le G_{\max},\quad \dots
\label{eq:constraints}
\end{equation}
and discard any candidate that violates a bound before it can be selected. Throughout we use
$\rho_{\min}=0.15$, variable coverage $\ge0.3$, $G_{\max}=0.6$ of the window and at least four observations per
sample. We additionally record a per-sample descriptor $r=g(\Oobs)\in\mathbb{R}^{8}$ of the realised process
(observation count and density, the mean, standard deviation and maximum inter-observation gap, temporal
coverage, variable coverage and recent density), computed from timestamps and masks alone and never from values,
which we use to report what a condition actually did to a batch.
